\chapter{\uppercase{Summary and Future Work}}
\label{chap:conclusion}

AgentOS provides a modular runtime for multi-agent applications in which planning, semantic agent selection, DAG execution, context management, scheduling, resource admission, sandboxed code execution, and frontend observability are treated as shared system services. The completed phase connects the gateway, planner, kernel, Redis event path, MongoDB task history, ChromaDB retrieval, scheduler, resource manager, agent runtime, and frontend activity interface into one execution pipeline.

The implementation is also bounded by explicit limitations. Scheduler dispatch is cooperative, token usage can be unavailable when a provider omits metadata, and restart recovery identifies interrupted work without restoring every active coroutine. The evaluation metrics defined in Appendix~\ref{app:metrics} provide the basis for measuring planning quality, execution reliability, scheduling fairness, resource compliance, context retrieval, event delivery, and code-runner safety; numerical results are reserved for controlled experiments.

Future work will therefore focus on durable restart resumption, true checkpoints inside long-running agents, complete capability enforcement, distributed scheduler workers, a hybrid local/cloud model router, cost-aware scheduling, formal workload benchmarks, and additional autonomous event triggers. These extensions can be added without changing the central separation between planner, scheduler, runtime, persistence, and observability layers.
